\chapter{Revisão Bibliográfica}
\label{cap:revisao}

\section{Microbioma e Metataxonômica}

O microbioma é definido como o conjunto de microrganismos — bactérias,
fungos, arqueas e vírus — que habitam um determinado ambiente,
juntamente com seus elementos genômicos e as condições
bióticas e abióticas que os envolvem \cite{berg2020microbiome}.
A metataxonômica, ramo da metagenômica, utiliza marcadores moleculares
conservados para identificar e quantificar esses microrganismos a
partir de amostras ambientais, sem necessidade de cultivo
\cite{knight2018best}.

Para bactérias, o marcador mais amplamente utilizado é o gene 16S
rRNA, cujas regiões hipervariáveis (V3--V4) permitem discriminar
gêneros e espécies com boa resolução taxonômica
\cite{johnson2019evaluation}. Para fungos, o espaçador interno
transcrito (ITS1/ITS2) da região ribossomal é o marcador recomendado
pela \textit{Fungal Barcoding Consortium}, pois apresenta variação
suficiente para separar espécies dentro de um mesmo gênero
\cite{schoch2012nuclear}.

\subsection{Pipeline QIIME2}

O \textit{Quantitative Insights Into Microbial Ecology 2} (QIIME2) é
atualmente o pipeline bioinformático mais adotado para análise
metataxonômica \cite{bolyen2019reproducible}. O fluxo de trabalho
compreende: importação dos arquivos FASTQ, controle de qualidade e
truncamento das leituras via DADA2 \cite{callahan2016dada2},
classificação taxonômica por Naive Bayes treinado em bases de referência
(SILVA 138 para 16S \cite{quast2012silva}; UNITE v10 para ITS
\cite{nilsson2019unite}), e exportação de tabelas de frequência OTU/ASV
no formato \texttt{phyloseq} \cite{mcmurdie2013phyloseq}.

O objeto \texttt{phyloseq} serializado (\texttt{.rds}) tornou-se o
formato padrão de entrada para a maioria das análises estatísticas
downstream implementadas na \textit{Rizoma}.

\subsection{Análises Estatísticas de Microbioma}

\subsubsection{ANCOM-BC2}

O ANCOM-BC2 (\textit{Analysis of Compositions of Microbiomes with
Bias Correction 2}) é um método robusto para detecção de abundâncias
diferenciais em dados de microbioma, projetado para corrigir o viés
de amostragem intrínseco a dados composicionais
\cite{lin2022analysis}. É o método adotado no projeto INOVAHERB para
comparar a composição do micobioma entre tratamentos fatoriais.

\subsubsection{MaAsLin2}

O MaAsLin2 (\textit{Multivariable Association with Linear Models 2})
modela associações entre features microbianas e metadados clínicos ou
ambientais por meio de modelos lineares generalizados, com suporte a
dados longitudinais e covariáveis múltiplas \cite{mallick2021multivariable}.
É empregado no projeto Pós-Fogo para identificar táxons cujas
abundâncias se correlacionam com o tempo decorrido após a queimada e
com variáveis edáficas mensuradas em campo.

\subsubsection{DESeq2}

O DESeq2 é um método bayesiano empírico originalmente desenvolvido
para RNA-Seq \cite{love2014moderated}, mas amplamente adaptado para
dados de microbioma após normalização por tamanho de biblioteca. Estima
a dispersão por gene e aplica \textit{shrinkage} nos log-fold changes
para maior estabilidade em estudos com poucas réplicas.

\subsubsection{SpiecEasi}

O SpiecEasi (\textit{Sparse InversE Covariance estimation for
Ecological Association Inference}) infere redes de co-ocorrência de
táxons a partir de estimativa esparsa de matrizes de covariância inversa
\cite{kurtz2015sparse}. Os nós da rede representam táxons e as arestas
representam correlações parciais, fornecendo uma representação mais
fiel das interações ecológicas do que correlações de Pearson ou
Spearman sobre dados composicionais.

\subsubsection{FUNGuild e PICRUSt2}

O FUNGuild é um banco de dados de anotação funcional para fungos,
acessível via API HTTP, que associa OTUs/ASVs a guildas ecológicas
(parasitas, simbiontes, saprófitas) \cite{nguyen2016funguild}. O
PICRUSt2 (\textit{Phylogenetic Investigation of Communities by
Reconstruction of Unobserved States 2}) realiza predição funcional
metagenômica a partir de dados 16S, estimando abundâncias de vias
metabólicas sem necessidade de sequenciamento metagenômico completo
\cite{douglas2020picrust2}.

\section{Arquiteturas para Plataformas Bioinformáticas}

\subsection{Microsserviços e Orquestração de Containers}

A abordagem de microsserviços decompõe uma aplicação em serviços
independentes, cada um com responsabilidade bem definida e
implantável de forma autônoma \cite{fowler2014microservices}. Para
plataformas científicas, essa arquitetura apresenta vantagens
significativas: permite escalar seletivamente os componentes mais
intensivos (como o worker de análises R), facilita a substituição de
implementações e isola falhas de modo a não comprometer o restante do
sistema.

O Kubernetes tornou-se o padrão de fato para orquestração de
containers em ambientes científicos \cite{kubernetes2022}. O k3s é
uma distribuição certificada e leve do Kubernetes, adequada a clusters
\textit{on-premises} com hardware limitado, como o servidor de cinco
nós AMD utilizado neste projeto.

\subsection{Filas de Jobs e Processamento Assíncrono}

A maioria das plataformas bioinformáticas utiliza sistemas de fila de
mensagens dedicados — como RabbitMQ, Apache Kafka ou Redis Streams —
para desacoplar a submissão de jobs de sua execução. Entretanto,
para cargas com dezenas a centenas de jobs diários, o PostgreSQL
oferece mecanismo equivalente via \texttt{LISTEN/NOTIFY} combinado
com \texttt{SELECT ... FOR UPDATE SKIP LOCKED}
\cite{brandl2021postgresql}, eliminando a necessidade de um serviço
adicional e simplificando a operação e a observabilidade.

\subsection{Armazenamento de Objetos e Dados}

O armazenamento de artefatos binários volumosos (arquivos FASTQ,
objetos R serializados, relatórios) requer atenção especial em
plataformas bioinformáticas. Uma alternativa ao uso de sistemas de
\textit{object storage} externos (como MinIO ou Amazon S3) é o
mecanismo nativo de \textit{Large Objects} do PostgreSQL, que permite
armazenar dados binários arbitrários referenciados por OIDs
(\textit{Object Identifiers}) e acessados via interface de
\textit{streaming} transacional \cite{postgresql2023docs}. Essa
abordagem centraliza backup, replicação e controle de acesso em um
único serviço, reduzindo a superfície operacional do cluster.

\subsection{CQRS e Event Sourcing em APIs Científicas}

O padrão \textit{Command Query Responsibility Segregation} (CQRS)
separa operações de escrita (comandos) das de leitura (consultas),
permitindo otimizar cada caminho de forma independente
\cite{young2010cqrs}. Em plataformas bioinformáticas, isso se traduz
em comandos transacionais gravados no PostgreSQL (enfileiramento de
jobs, confirmação de upload) e consultas de leitura rápida servidas
pelo Elasticsearch, que mantém índices otimizados para agregações e
buscas textuais sobre resultados analíticos.

\section{Visualização de Dados de Microbioma}

\subsection{Análise de Coordenadas Principais (PCoA)}

A PCoA é a técnica de ordenação mais empregada para representar
visualmente a diversidade beta de comunidades microbianas
\cite{legendre2012numerical}. A partir de uma matriz de distâncias
(Bray-Curtis, UniFrac, Jaccard), a PCoA projeta as amostras em um
espaço de menor dimensão que maximiza a variância explicada, permitindo
identificar agrupamentos por tratamento, tempo ou outra variável
experimental. Neste trabalho, seis PCoAs — duas por projeto — compõem
o painel central do TCC.

\subsection{Redes Microbianas}

Redes de co-ocorrência microbiana são representadas como grafos nos
quais nós correspondem a táxons e arestas a correlações significativas
entre seus perfis de abundância \cite{faust2012microbial}. A
identificação de táxons \textit{hub} — nós com alto grau de
conectividade, também denominados táxons-chave (\textit{keystone taxa})
— é relevante para compreender a estabilidade e a resiliência do
ecossistema microbiano \cite{berry2014deciphering}.

\section{Tecnologias de Frontend para Ciência}

Plotly.js é uma biblioteca de visualização interativa amplamente
adotada em aplicações científicas por seu suporte nativo a gráficos
2D e 3D, exportação em alta resolução e integração com React via
\texttt{react-plotly.js} \cite{plotly2015}. Para redes complexas,
Cytoscape.js oferece renderização eficiente de grafos com centenas de
nós e arestas, com suporte a layouts automáticos e interatividade
\cite{franz2016cytoscape}.

% ─── TODO: ampliar conforme evolução do texto ────────────────────────────────
% Sugestões de seções adicionais:
%   \section{Reprodutibilidade em Bioinformática}
%   \section{Trabalhos Relacionados}
